% ED26-BM-OHTR2012-6: full Morrison-Knowlton2012 methods treatment.
\chapter{Neurocognitive Methods in Higher Cognition}
\label{sec:p19-neurocognitive-methods}

Morrison and Knowlton distinguish mapping brain--behavior relationships
from using neural evidence to constrain competing cognitive theories.
Several mechanisms can reproduce a behavioral result; an additional
neural prediction can exclude some, provided the measurement and
the proposed brain/process bridge are defensible. Their chapter
surveys neuropsychology, structural and functional imaging, electrical
recordings, magnetic perturbation and computational models, emphasizing
prefrontal cognition and cooperative networks rather than a single
reasoning center \cite[chap.~6, pp.~67--89]{MorrisonKnowlton2012Methods}.

The source is a 2012 methods presentation. Its equipment ranges,
availability judgments and forecasts are dated descriptions, not
current clinical protocols. Nothing here recommends diagnosis,
treatment, stimulation or participant eligibility. Its own note~1
distinguishes \emph{cognitive neuropsychology}, the use of patient
dissociations to investigate cognition, from clinical neuropsychology's
psychometric appraisal. Note~2's diagnostic terminology supplies
the historical grouping used in the studies, not instructions for
classifying a person.

\section{What a neural observation can constrain}
\label{sec:p19-ncm-inference}

\begin{hbdefinition}{def:p19-ncm-claims}{Different neural and cognitive claims}
A \emph{correlate} covaries with a task, response or other measure.
\emph{Localization} assigns that measured signal or relevant contribution
to a specified anatomical site under a spatial model. \emph{Necessity}
claims that the function would fail without the relevant contribution
under the stated conditions. A \emph{perturbation effect} compares
outcomes under an intervention and its appropriate controls.
\emph{Mechanistic identification} additionally specifies the information
and operations by which the contribution supports behavior.

These are not synonyms. A region activated during a task need not
be necessary for its defining process; a necessary contribution can
be sensory, linguistic, motivational or motor rather than the
reasoning operation of interest. An intervention's effect also
does not establish that one local circuit is sufficient for the
whole behavior.
\end{hbdefinition}

The anatomical orientation in Figure~6.1 distinguishes frontal,
parietal, temporal and occipital lobes and narrower regions such
as dorsolateral and rostrolateral prefrontal cortex, DLPFC/RLPFC,
inferior frontal gyrus, IFG, and anterior temporal cortex, ATC.
Labels based on anatomy, historical investigators or functional
associations have different origins. The drawing locates names;
it does not demonstrate an exclusive function for every region.

Reverse inference moves from observed activity to the claim that
a particular cognitive process occurred. Forward evidence that
a process often activates a region is not enough: other processes
may activate it too, and the target process's prevalence in the
actual comparison matters. Baselines and task phases must also
be specified. A language-associated region active while premises
are read does not by itself show that subsequent deduction uses
a linguistic rule representation. This methodological distinction
develops the source's repeated warning about incidental processes,
not a clinical interpretation of a participant's scan.

\section{Patient cases and dissociation assumptions}
\label{sec:p19-ncm-patients}

The source contrasts historical localizationism with an aggregate
field account and recounts Broca's patients Leborgne and Lelong.
Modern inspection showed more extensive damage than the original
local description suggested. Similar presentations can also include
grammatical comprehension deficits, motivating an agrammatism
interpretation rather than a simple production/comprehension split.
The lesson concerns changing behavioral characterization and lesion
mapping, not diagnosing aphasia from one task
\cite[pp.~68, 72 and figs.~6.1--6.2]{MorrisonKnowlton2012Methods}.

\begin{hbdefinition}{def:p19-ncm-dissociation}{Single and double dissociation as evidence patterns}
A single dissociation has impaired performance on one function
with another relatively spared. A double dissociation adds the
complementary profile: another case or group is relatively impaired
on the second while sparing the first. Figure~6.2 illustrates
the two profiles. Two cases with the same production-impaired
profile do not alone establish a double dissociation.

Interpreting the pattern requires comparable task sensitivity,
appropriate control performance and a model of the effects of
damage. Unequal task difficulty, shared resource demands, disconnection,
compensation and different strategies can complicate the inference.
An observed dissociation constrains a theory; it does not automatically
prove anatomically isolated independent modules.
\end{hbdefinition}

Skill learning with poor memory for practice episodes motivates
the source's declarative/procedural-memory contrast. For reasoning,
its frontal/temporal patient comparison separates relational integration,
interference control and access to semantic knowledge. Single-relation
matrix tasks can be spared while tasks combining changes across
rows and columns are impaired in the frontal group. Temporal-group
semantic deficits can instead affect familiar materials while
arbitrary-content reasoning remains relatively possible. These are
chapter reports attributed to Waltz, Morrison, Krawczyk and related
studies, not prevalence estimates for all neurological populations
\cite[pp.~68, 72--73]{MorrisonKnowlton2012Methods}.

Damage location and extent vary; focal unilateral lesions may
permit compensation, and degenerative groups are neither random
lesions nor homogeneous circuit removals. Structural volume--behavior
associations help characterize these differences but do not themselves
make a damaged site necessary for one unique computation. The
source's semantic/arbitrary-circuit interpretation therefore remains
an explanatory proposal requiring task-sensitive comparisons.

\section{Measurement families and the source's comparative table}
\label{sec:p19-ncm-measurements}

Table~6.1 compares thirteen methods across lag, acquisition timing,
spatial scale, advantages and constraints. Those fields answer
different questions. A millisecond sampling interval is not a
millisecond hemodynamic response, and it is not a rate expressed
in hertz. Anatomical image resolution is not identical to the
resolution of functional source inference.

\begin{center}
\begin{tabular}{@{}p{0.27\textwidth}p{0.29\textwidth}p{0.36\textwidth}@{}}
\toprule
Source method & Recorded information & Main inferential limitation\\
\midrule
Patient neuropsychology & Behavioral deficit/spared-function profiles,
with damage characterization & Nonrandom heterogeneous damage,
compensation and task sensitivity.\\
CT & Structural X-ray reconstruction & Anatomy does not directly
measure an ongoing cognitive operation.\\
Structural MRI/VBM & Tissue contrast, thickness and normalized regional
volume comparisons & Segmentation, registration and group comparisons
require models; volume is not a synapse count.\\
DTI & Directional water-diffusion measurements & Tissue/connectivity
interpretation is indirect, not direct observation of messages.\\
Single/multi-unit recording & Local electrode activity & Invasive,
restricted sites and species/task transfer.\\
ECoG & Electrical activity at selected cortical sites & Selected
participant/site coverage, volume conduction and pathology.\\
EEG/ERP/ERO & Scalp voltages, time-locked means or frequency analyses
& Excellent timing does not uniquely locate a generator or identify a process.\\
PET & Tracer-related metabolic/receptor measurements & Indirect neural
interpretation, coarse timing and invasive tracer procedures.\\
fMRI & BOLD hemodynamic measurements & Delayed vascular signal,
baseline/analysis dependence and no necessity proof from activation.\\
MEG & Magnetic fields associated with neural currents & Source-model
ambiguity remains despite different spatial sensitivity from EEG.\\
fNIRS/EROS & Hemoglobin-related optical measures/fast optical scattering
& Different physiological targets; superficial coverage and model-dependent inference.\\
TMS & Perturbation and subsequent behavior/recording & Controls,
timing, field spread and network consequences limit localization.\\
\bottomrule
\end{tabular}
\end{center}

This is an editorial reconstruction of the comparison, not a new
benchmark. The original table's spatial orders range from
10 micrometers--1 millimeter for unit recording through roughly
1--10 millimeters for ECoG, optical methods, TMS and lesion characterization,
to roughly 10--100 millimeters for EEG/MEG, explicitly questioned
there. Structural/fMRI entries have broad millimeter--decimeter
ranges; PET is listed at centimeter--decimeter scale. They are
2012 approximate table entries, not guarantees of cognitive localization.
Its fMRI acquisition intervals are 250 milliseconds--1 second,
with 1--2-second onset lag and a response developing over seconds;
fNIRS also has a hemodynamic lag despite fine sampling. The EROS
entry lists about 100 milliseconds of lag. The table's question
marks and unavailable timing fields remain uncertainty, not zero
\cite[table~6.1, pp.~69--71]{MorrisonKnowlton2012Methods}.

\subsection{Structural methods}

CT and then MRI made damage accessible without waiting for a
postmortem examination. MRI supports tissue contrast, volumetry
and cortical-thickness analysis; VBM introduces spatial normalization
and voxelwise comparisons. The source's developmental reasoning
example associates adolescent prefrontal gray-matter reductions
with behavior. It suggests synaptic/connectivity change, but MRI
volume alone does not demonstrate pruning or prove its causal
effect on reasoning. DTI measures diffusion characteristics that
can inform white-matter models; it is not a direct assay of
fiber integrity or information flow. These distinguish measured
quantities from biological interpretations
\cite[pp.~73--74]{MorrisonKnowlton2012Methods}.

\subsection{Electrical measurements and working-memory control}

Local unit recordings can constrain when activity occurs at a
sampled site. The source's macaque delayed-match example separates
sample presentation, maintenance delay and response. Its cortical-
cooling comparison supports an interaction account in which
prefrontal activity helps maintain posterior representations, rather
than a simple isolated storage buffer. Species, task and intervention
conditions constrain transfer to human reasoning
\cite[pp.~74--75]{MorrisonKnowlton2012Methods}.

ECoG provides selected human cortical recordings; the chapter's
examples come from clinically obtained access rather than a
random sample of intact brains. Cross-frequency coupling motivates
communication/binding hypotheses, but correlated phase and power
do not count neurons or uniquely establish what content is bound.
Scalp EEG is a mixed electrical signal, not separate spike trains
from identified neurons. Its spatial inverse problem permits
multiple source configurations producing similar recorded voltages.

ERPs average measurements time-locked to a stimulus or response.
Figure~6.3 labels illustrative positive/negative peaks P1, N1, P2, N2, P3;
the sign convention, reference, latency and task must be known before
a component is interpreted. Averaging can reduce non-time-locked
noise under suitable assumptions, but systematic artifacts and
trial-to-trial timing variation do not disappear by definition.
Event-related oscillation analyses retain information about frequency,
phase and time that an average waveform can obscure. A Fourier
or wavelet transformation does not directly estimate the number
of neurons participating in each cognitive process.

The source's P30 auditory-click patient comparison suggests sensory
gating by DLPFC. Its insight examples describe alpha activity
before a solution and gamma activity near the reported answer;
the prepared-mind example concerns prediction before problem
presentation. These are different contrasts, not an identical
neural signature of insight. Except for the separately acquired
Vogel anchor below, the chapter does not supply each original
sample, trial denominator, effect interval or complete analysis.
Interpretations as gating or semantic integration remain competing-
mechanism claims, not uniquely identified processes
\cite[pp.~75--77]{MorrisonKnowlton2012Methods}.

\subsection{The actual Vogel 2005 anchor}

Vogel and colleagues studied 15 neurologically normal college
students aged 19--28 in each of three experiments, with 200--240
trials per condition per participant. These are participants and
within-person trials, not hundreds of independent people. A
visual-memory estimate supplied a median split into
high/low groups; group membership was not randomized
\cite[pp.~500--503]{VogelEtAl2005Selection}.

Experiment~1 used a cued left/right hemifield. Participants remembered
red rectangles' orientations and judged whether a test array one
second later changed. Conditions were two red items, four red items,
or two red plus two irrelevant blue items. The EEG outcome was
\emph{contralateral delay activity}, CDA, from posterior sites;
the original paper measures a lateralized difference during
300--900 milliseconds after array onset. It is not the contingent
negative variation, CNV, as named in the 2012 survey. Twenty-two
electrode sites were recorded, with blink/large-eye-movement trial
rejection. Missing trials and retained averages therefore matter.

Distractor-present CDA was nearer the two-item condition in the
high group and the four-item condition in the low group. The
filtering-score/capacity correlation was \(r=0.69\), reported
\(p<0.001\); no confidence interval is given. Experiment~2
used easier location selection and reported \(r=0.62\),
\(p<0.001\). Experiment~3 introduced a second array after
500 milliseconds and distinguished appending red items from
excluding green ones; its reported relation was \(r=0.57\),
\(p<0.01\). Appending was relatively preserved across groups
while exclusion differed. Changed task versions support a bounded
generalization of the filtering interpretation, not a universal
capacity theory \cite[pp.~500--502]{VogelEtAl2005Selection}.

\begin{hbdefinition}{def:p19-ncm-scores}{The primary's capacity and filtering measures}
For the capacity estimate the paper uses \(K=S(H-F)\), where
\(S\) is array size, \(H\) the hit proportion and \(F\)
the false-alarm proportion. This is a supplied task model and
estimator, not a direct count of an individual's neural storage slots.

In the separate filtering-score context the source reuses \(F\)
for four-item mean CDA amplitude, \(D\) for distractor-present
amplitude and \(T\) for two-item amplitude. Its published score is
\[
                       \alpha=\frac{F-D}{F-T}.
\]
Here those symbols are amplitudes, not response proportions.
The score measures the distractor condition's relative position
between two physiological comparison conditions.
\end{hbdefinition}

\begin{hbtheorem}{thm:p19-ncm-score}{proved=annex:p19.ncm-score}{The actual filtering-score range and boundary}
For real amplitudes with \(F\ne T\), the score equals one
when \(D=T\), zero when \(D=F\), and lies in \([0, 1]\)
exactly when \(D\) lies between \(T\) and \(F\), including
endpoints. If \(F=T\), the ratio is undefined. For the separate
capacity formula with \(S>0\) and \(H, F\in[0, 1]\),
\(-S\leq K\leq S\), and nonnegative \(K\) additionally
requires \(H\geq F\). Neither formula by itself establishes
that the reported score equals an intrinsic cognitive capacity.
\end{hbtheorem}

The primary explicitly allows filtering scores outside zero--one
mathematically while reporting none in its observations. Its
interpretation is that capacity differences need not \emph{simply}
reflect storage space: allocation efficiency can contribute.
The 2012 survey's stronger claim that the groups do not differ
in capacity is not adopted. Nor does a posterior CDA association
localize its causal source uniquely to PFC. The primary proposes
a prefrontal bias signal as plausible, not as an intervention result
\cite[p.~502]{VogelEtAl2005Selection}.

\section{Hemodynamics, optics and network modeling}
\label{sec:p19-ncm-hemodynamics}

PET uses tracer-related measurements and fMRI uses blood-oxygenation-
dependent, BOLD, signal. Both require a physiological bridge from
neural activity to the measured quantity. Tracer choice can address
different metabolic/receptor questions; it does not make PET a
direct measure of every neurotransmitter's firing. Figure~6.4
illustrates transient oxygen demand followed by a vascular response.
The increased positive-BOLD phase has reduced relative
deoxyhemoglobin following that response. The displayed curves
are explanatory schematics, not
new participant data \cite[pp.~77--78]{MorrisonKnowlton2012Methods}.

Event-related fMRI can separate modeled phases and intermix trials,
but blood dynamics blur timing. Removing error trials changes
the estimand; it does not reconstruct an unobserved error-free
mind. Activity varies with logical form, meaningful content, belief
conflict and perceptual demands. The source contrasts arbitrary
conditional problems, meaningful bicycle/weather problems and
transitive spatial relations. Its conditional examples are not
categorical syllogisms merely because the survey uses that label.
Every inference benchmark requires its actual premise interpretation.

The survey does not find unequivocal imaging support for mental
logic over mental models or vice versa. The premise that a
rule theory must use linguistic cortex and a model theory right
visuospatial cortex is itself a bridge hypothesis. Separating
premise reading from solution phases can constrain it, but a
region's temporal contrast does not prove a unique representational
format. The claim is about possible separation of language
processing and reasoning, not reasoning being independent of thought
\cite[pp.~78--79]{MorrisonKnowlton2012Methods}.

The Sonty language-network example contrasts semantic with letter
matching in patients and controls. Poorer behavior occurred without
simply reduced local activation; a fitted dynamic-causal-modeling,
DCM, connectivity difference supplied another account. The modeled
coupling/performance association does not prove that reduced
communication uniquely caused the deficit. DCM estimates depend
on a generative model, selected regions, connections and comparison;
it is not direct observation of individual synapses. Functional
connectivity based on correlated BOLD is likewise distinct from
directional causal influence or an anatomical tract
\cite[pp.~80, 84]{MorrisonKnowlton2012Methods}.

MEG measures magnetic fields associated with neural currents.
Its different field/source sensitivities can constrain localization
but do not remove the inverse problem. The source's child/adult
n-back example permits equal behavioral performance with different
signals; neither identical behavior nor different signals alone
settles which complete mechanism is used.

fNIRS estimates changes involving oxy/deoxyhemoglobin, with rapid
acquisition but delayed hemodynamics. EROS targets faster optical
scattering rather than that same blood response. Both have
superficial-coverage limitations in the source. The chapter's
48-participant fNIRS belief-bias example crosses short/long time
with congruent/incongruent problems and associates right IFG
activity with some accuracy contrasts. Congruence here should
separate logical validity from conclusion believability; factual
truth alone is not logical validity. The source supplies a sample
count and directional survey but no full original effect/interval/
design record. An activity correlation does not establish the
claimed inhibition mechanism or necessity
\cite[pp.~80--81]{MorrisonKnowlton2012Methods}.

\section{TMS: intervention with limited localization}
\label{sec:p19-ncm-tms}

TMS changes electromagnetic conditions rather than only recording
a signal. The relevant evidence compares behavior under the
intervention with suitable conditions controlling site, timing,
sound, sensation, expectation and task demands. Perturbation may
alter linked networks or other processes required to perform
the task. Enhancement and disruption are not interchangeable
with permanent focal damage, or with sufficient activation of
one reasoning module.

The source's right-IFG follow-up reports increased belief-bias
under repetitive TMS on incongruent trials. This supports a
contribution under the studied task/control conditions, not a
deduction that IFG contains only a belief-inhibition operation.
The chapter does not supply the complete original sample,
stimulation/control/effect/uncertainty record. It also notes
surface accessibility and the difficulty of testing cortical--
subcortical interactions this way. The treatment discussion and
categorical eligibility statements in its dated table are not
adopted as current clinical guidance or an intervention recommendation
\cite[pp.~81--82]{MorrisonKnowlton2012Methods}.

\section{A worked prefrontal/semantic comparison}
\label{sec:p19-ncm-analogy}

The separately acquired Morrison 2004 primary studied picture and
verbal analogies in frontal/temporal FTLD groups and age/education-
matched controls. Experiment~1 included 7 frontal patients,
6 temporal patients and 8 controls; experiment~2 included 7, 6
and 10 respectively. These are the actual task-specific samples,
not randomly assigned lesions or a representative population.
Selection included abnormal imaging and the source's MMSE
threshold. That characterizes recruitment; it is not a diagnostic
rule offered here \cite[pp.~267--268]{MorrisonEtAl2004Analogy}.

Ten picture pairs permitted a salient feature match or a relational
match. Participants inspected each for ten seconds, then indicated
what went with a pointed-out object; there were no practice
examples and one predetermined order. Both patient groups made
fewer relational and more featural choices than controls. Frontal
relational choices more often benefited from perceptual alignment.
That pattern leaves working-memory integration, inhibition and
semantic access as competing contributors, not one uniquely located
faculty \cite[pp.~262--264, 267]{MorrisonEtAl2004Analogy}.

Experiment~2 used 24 two-choice verbal analogies selected from
200 candidates with association/difficulty norms from separate
UCLA undergraduate samples, 150 for association and 54 for
difficulty. A semantic-facilitation index, SFI, contrasted normalized
association of the correct versus distractor pair. Positive,
neutral and negative conditions used thresholds around plus/minus
0.3. The source's play:game::give:? example pits \emph{party}
against the more strongly associated \emph{take}. Random problem
order followed five practice items. Outcome was answer accuracy,
not a direct neural measure or proof of the association's causal origin.

Frontal-group accuracy was selectively poorer when association
favored the distractor; temporal-group accuracy was more uniformly
reduced. The source reports the frontal differential contrast
\(F(1, 19)=5.0\), \(p=0.04\), while the analogous temporal
contrast was not detected. Figure~4 uses one-standard-error bars;
no confidence interval or participant-level reanalysis is supplied.
The reported analysis unit/denominator must not be inferred solely
from the degrees of freedom or trial count. Small clinical groups,
heterogeneous damage, item selection and fixed-order picture testing
limit generalization \cite[pp.~265--268]{MorrisonEtAl2004Analogy}.

Its LISA simulation used 40 trials per problem type and changed
mapping learning/inhibition to model frontal deficits versus semantic
connections for temporal deficits. It reproduced qualitative patterns,
with tested parameter assumptions rather than measured neuronal
values. Neither frontal manipulation alone reproduced the selected
pattern in that model. This is a useful competing-mechanism
boundary, not independent proof that actual lesions change precisely
those parameters or that no other architecture can explain the
results \cite[pp.~266, 268--270]{MorrisonEtAl2004Analogy}.

\section{Neurocomputational bridges and the prefrontal alternatives}
\label{sec:p19-ncm-models}

ACT-R's production architecture and its task-specific assumptions
can predict behavior; associating its representation, retrieval
and response functions with regions introduces an additional
bridge. The chapter's Tower-of-Hanoi example is an external
model/imaging account, not an acquired executable ACT-R instance.
Behavioral fit alone is not neural validation.

Figure~6.5 develops LISA's proposed bridge in five panels: distributed
semantic representations and conjunctive long-term codes; dynamic
working-memory proxy/role-binding units; candidate retrieval based
on integrated activation; synchronized driver/recipient mapping
through rapidly learned links; and increased anterior-prefrontal
involvement for higher-order relations. Its lover/loved example
binds the two roles in distinct temporal phases rather than mixing
the object features into an unordered list. DLPFC maintenance,
IFG competition control, ATC/hippocampal representation and
RLPFC mapping assignments are hypotheses. Observed rhythmic
activity alone does not verify each proxy unit, Hebbian connection
or synchrony-based relational binding
\cite[pp.~82--84 and fig.~6.5]{MorrisonKnowlton2012Methods}.

The source distinguishes crystallized knowledge from fluid reasoning
and surveys task-complexity-dependent recruitment. It does not
turn a scan into an intelligence diagnosis. Common lateral
prefrontal/anterior-cingulate recruitment can reflect a shared
control operation or flexible support of several operations.
The hierarchical proposal distinguishes sensory-response, context
and episodic-memory demands; medial/orbital social/emotional
associations have other comparisons. RLPFC integration and
abstraction accounts remain explicit alternatives. A region's
name or greater activation does not decide between them
\cite[pp.~84--85]{MorrisonKnowlton2012Methods}.

The chapter's three future priorities are cortical connectivity,
integration of spatial/temporal measurement, and neurally plausible
computational hypotheses. Convergence of methods is valuable when
their independent assumptions and predicted contrasts constrain
one another. Repeating one localization assumption in several
analyses is not independent corroboration.

\section{Consultation and related evidence interfaces}
\label{sec:p19-ncm-use}

Goel's Chapter~\ref{sec:p19-rationality-1-4} offers a particular
reasoning-systems account, not the present chapter's whole methods
survey. Knauff's Chapter~\ref{sec:p19-rationality-13-3} separates
visual detail, spatial organization and task performance; an occipital
signal need not identify pictorial reasoning. Cognitive representation
in Chapter~\ref{sec:p19-cognitive-knowledge-representation} supplies
encoding/process distinctions that a neural bridge must preserve.
These are complementary interfaces, not duplicate source credits.

\begin{hbmethod}{def:p19-ncm-use}{Match the neural method to the inference actually needed}
Specify the population, task, behavioral measure and proposed mechanism.
Name the neural observable, physiological bridge, baseline, time phase,
spatial/source model and processing steps. Distinguish acquisition
timing from physiological lag and analysis resolution.

For patient evidence characterize lesion/degeneration and spared
performance, including sensitivity and compensation alternatives.
For activation evidence test selectivity and competing processes
before reversing the inference. For perturbation state intervention,
controls, network effects and the task component actually changed.
For connectivity specify observed correlation versus model-based
directional coupling and what could distinguish other models.

Construct the task-specific evidence record in
Method~\ref{def:p19-evidence-record}. Keep participants, trials,
conditions and model runs as distinct denominators. A nonsignificant
difference is not proof of equality; a successful fitted simulation
is not a unique brain mechanism. State transfer and replication
limits and preserve unresolved source-study boundaries.
\end{hbmethod}

The useful success case combines retained behavioral contrasts
with method-specific neural predictions and an actually controlled
perturbation or dissociation. Failure cases include reading capacity
alone from a mixed ERP, inferring a specific thought from generic
PFC activation, calling correlated BOLD a proved causal pathway,
or transferring one patient task to all reasoning. The source
provides tools for distinguishing these questions, not a universal
neural verdict of rationality.

\index{neurocognitive methods}
\index{neural evidence!reverse inference}
\index{neuropsychology!single and double dissociation}
\index{EEG!ERP and oscillations}
\index{fMRI!BOLD and inference}
\index{TMS!perturbation and localization}
\index{prefrontal cortex!competing mechanisms}
